\phantomsection
\addcontentsline{toc}{chapter}{课程信息与原教学进度}
\markboth{课程信息与原教学进度}{}
\noindent{\color{structurecolor}\bfseries\fontsize{14.4}{20}\selectfont 课程信息与原教学进度}\par

原课程为 \textit{Probability and Random Variables}, MIT 18.600, 2019年秋季, 授课教师 Scott Sheffield, 开课单位为麻省理工学院数学系, 本科层次. 官方先修课为18.02多变量微积分. 2015年春季以前的课号为18.440; 改号后的18.600表示18.6xx系列的基础课程, 本书所收旧卷仍属于同一课程的历年训练材料.

\noindent\textbf{课时与学习安排.} 每周3次讲授, 每次1小时. 官方教师访谈列出39次讲授、13次助教习题课, 习题课每周1小时, 每周总学习投入约12小时, 包括课堂以外的作业与复习. 官方Calendar另把期末考试列为第40次活动; 未公布期末考试的课时, 本书不据此把40次活动都算成40小时讲授.

\noindent\textbf{内容大纲.} 计数与概率公理, 条件概率和贝叶斯公式, 独立性, 离散和连续随机变量, 分布函数, 二项、几何、超几何及Poisson分布, 均匀、指数、正态、Gamma和Beta分布, 联合分布, 期望、方差与协方差, 条件期望和母函数, Chebyshev不等式、大数定律与中心极限定理. 原教学进度还包括有限马尔可夫链、熵、鞅、可选停止以及风险中性定价.

\noindent\textbf{教材与阅读.} 必修教材为Sheldon Ross, \textit{A First Course in Probability}, 第8版, Pearson Prentice Hall, 2009, ISBN 9780136033134. 下表中的节号均指该版. 官方还推荐Charles Grinstead和J. Laurie Snell的\textit{Introduction to Probability}, 并在\href{https://ocw.mit.edu/courses/18-600-probability-and-random-variables-fall-2019/pages/readings/}{Readings目录}提供概率史、不可测集合、大数定律及金融模型等补充入口. 本书不复制这些外部教材或文章, 也不把其题号当作已取得的题面.

\noindent\textbf{作业与考核.} 全学期10份作业, 考试周不另布置作业. 成绩组成是作业20\%, 两次期中考试合计40\%, 期末考试40\%. 官方考试说明为闭卷且不带公式纸; story sheet是备考材料. 原目录未给出两次期中各自占比, 不从合计40\%推定各占20\%.

\noindent\textbf{资源入口.} \href{https://ocw.mit.edu/courses/18-600-probability-and-random-variables-fall-2019/pages/syllabus/}{Syllabus}, \href{https://ocw.mit.edu/courses/18-600-probability-and-random-variables-fall-2019/pages/calendar/}{Calendar}, \href{https://ocw.mit.edu/courses/18-600-probability-and-random-variables-fall-2019/pages/assignments/}{Assignments}, \href{https://ocw.mit.edu/courses/18-600-probability-and-random-variables-fall-2019/pages/exams/}{Exams}及\href{https://ocw.mit.edu/courses/18-600-probability-and-random-variables-fall-2019/pages/instructor-insights/}{Instructor Insights}均为MIT官方页面. 该课程页面没有提供课堂视频系列; 其他媒体、访谈及补充链接作为阅读入口, 未作逐字媒体转录. 本表保留原课程顺序, 本书正文则按数学依赖重组.

\begin{longtable}{@{}p{.07\linewidth}p{.44\linewidth}p{.20\linewidth}p{.19\linewidth}@{}}
\toprule
讲次 & 原课主题 & Ross阅读节号 & 课程节点\\\midrule
\endfirsthead
\toprule 讲次 & 原课主题 & Ross阅读节号 & 课程节点\\\midrule
\endhead
1 & 排列与组合 & 1.1--1.3 & \\
2 & 多项式系数与计数 & 1.4--1.5 & \\
3 & 样本空间与集合 & 2.1--2.2 & \\
4 & 概率公理 & 2.3--2.4 & \\
5 & 概率与等可能性 & 2.5--2.7 & 作业1截止\\
6 & 条件概率 & 3.1--3.2 & \\
7 & 贝叶斯公式与独立事件 & 3.3--3.5 & \\
8 & 离散随机变量 & 4.1--4.2 & 作业2截止\\
9 & 离散随机变量的期望 & 4.3--4.4 & \\
10 & 方差 & 4.5 & 作业3截止\\
11 & 二项变量、重复试验与投资组合 & 4.6 & \\
12 & Poisson随机变量 & 4.7 & \\
13 & Poisson过程 & 9.1 & 作业4截止\\
14 & 其他离散随机变量 & 4.8--4.9 & \\
15 & 第一次期中复习 & 无指定阅读 & \\
16 & 第一次期中考试 & 无指定阅读 & 无讲义\\
17 & 连续随机变量 & 5.1--5.3 & \\
18 & 正态随机变量 & 5.4 & 作业5截止\\
19 & 指数随机变量 & 5.5 & \\
20 & 其他连续随机变量 & 5.6--5.7 & \\
21 & 联合分布函数 & 6.1--6.2 & 作业6截止\\
22 & 独立随机变量之和 & 6.3--6.5 & \\
23 & 和的期望 & 7.1--7.2 & \\
24 & 协方差与条件期望习题 & 7.3--7.4 & 作业7截止\\
25 & 条件期望 & 7.5--7.6 & \\
26 & 矩母函数 & 7.7--7.8 & \\
27 & 弱大数定律 & 8.1--8.2 & 作业8截止\\
28 & 第二次期中复习 & 无指定阅读 & \\
29 & 第二次期中考试 & 无指定阅读 & 无讲义\\
30 & 中心极限定理 & 8.3 & \\
31 & 强大数定律与Jensen不等式 & 8.4--8.5 & \\
32 & 马尔可夫链 & 9.2 & 作业9截止\\
33 & 熵 & 9.3--9.4 & \\
34 & 鞅与可选停止定理 & 官方补充链接 & \\
35 & 鞅与风险中性概率 & 官方补充讲义 & \\
36 & 风险中性概率与Black--Scholes & 官方补充链接 & \\
37 & 期末复习 & 无指定阅读 & 作业10截止\\
38 & 期末复习续 & 无指定阅读 & \\
39 & 期末复习续 & 无指定阅读 & \\
40 & 期末考试 & 无指定阅读 & \\
\bottomrule
\end{longtable}
